\documentclass{beamer}
\usepackage[T1]{fontenc}
\usepackage[utf8]{inputenc}
\usepackage[brazilian]{babel}
\usepackage{url}

\usetheme{Szeged}

\title{Qualidade de Serviço (QoS)}
\subtitle{Redes de Computadores}
\author[Marco, DKM]{Marco Henry\inst{1} \and Danilo Kotaka\inst{2}}
\institute[UEL]{Universidade Estadual De Londrina}
\date{\today}

\AtBeginSection[]
{
  \begin{frame}
    \frametitle{Sumário}
    \tableofcontents[currentsection]
  \end{frame}
}

\begin{document}

  \frame{\titlepage}

  \begin{frame}
    \frametitle{Sumário}
    \tableofcontents
  \end{frame}

  % =============================================
  \section{Introdução}
  % =============================================

  \begin{frame}
    \frametitle{Contexto Histórico}
    \begin{itemize}
      \item<1-> Redes separadas: voz em circuitos dedicados, dados em redes distintas
      \item<2-> Convergência: voz, vídeo e dados na mesma infraestrutura
      \item<3-> Apps interativas (VoIP, streaming, videoconferência) exigem qualidade garantida
      \item<4-> IoT amplia a demanda: sensores industriais precisam de comunicação em tempo real
    \end{itemize}
  \end{frame}

  % =============================================
  \section{Definição de QoS}
  % =============================================

  \begin{frame}
    \frametitle{O que é QoS?}
    \begin{block}{Definição}
      \textbf{Qualidade de Serviço (QoS)} é o uso de mecanismos e tecnologias
      que atuam em uma rede para controlar o tráfego e assegurar a boa
      \textit{performance} de aplicações críticas com capacidade de rede
      limitada.
    \end{block}
    \vspace{0.5em}
    \begin{itemize}
      \item<2-> Marca pacotes para identificar tipos de serviço
      \item<3-> Cria filas virtuais separadas por prioridade
      \item<4-> Reserva largura de banda para aplicações críticas
    \end{itemize}
  \end{frame}

  \begin{frame}
    \frametitle{Serviços que Requerem QoS}
    \begin{columns}
      \column{0.5\textwidth}
      \begin{itemize}
        \item IPTV
        \item \textit{Online gaming}
        \item \textit{Streaming} de mídia
      \end{itemize}
      \column{0.5\textwidth}
      \begin{itemize}
        \item Videoconferência
        \item Video on Demand (VOD)
        \item Voice over IP (VoIP)
      \end{itemize}
    \end{columns}
    \vspace{1em}
    \begin{alertblock}{Por quê?}
      Esses serviços são \textbf{inelásticos}: têm requisitos mínimos de largura
      de banda e são altamente sensíveis a atraso e jitter.
    \end{alertblock}
  \end{frame}

  % =============================================
  \section{Elementos Avaliados}
  % =============================================

  \begin{frame}
    \frametitle{Largura de Banda e Atraso}
    \begin{block}{Largura de Banda (\textit{Bandwidth})}
      Capacidade máxima do enlace. QoS aloca banda por fila de tráfego.
    \end{block}
    \vspace{0.5em}
    \begin{block}{Atraso (\textit{Delay} / Latência)}
      Tempo de entrega origem-destino. Afetado por \textit{queuing delay} em
      congestionamentos.
      \begin{itemize}
        \item Em VoIP: latência alta causa sobreposição de áudio
        \item QoS cria filas de prioridade para reduzir atraso
      \end{itemize}
    \end{block}
  \end{frame}

  \begin{frame}
    \frametitle{Variação de Atraso (Jitter)}
    \begin{block}{Jitter}
      Irregularidade na velocidade de chegada dos pacotes. Causa quebras em
      áudio e vídeo.
    \end{block}
    \vspace{0.5em}
    \begin{block}{Buffer de Jitter (\textit{Playout Buffer})}
      Armazena pacotes temporariamente e os entrega em intervalos uniformes.
      \begin{itemize}
        \item<2-> \textbf{Estático}: tamanho fixo, configurado manualmente
        \item<3-> \textbf{Adaptativo}: ajusta tamanho conforme condições da rede
        \item<4-> Trade-off: adiciona latência constante, elimina variação
        \item<5-> VoIP: tipicamente 50--150 ms de buffer de áudio
      \end{itemize}
    \end{block}
  \end{frame}

  \begin{frame}
    \frametitle{Perda de Pacotes e Outras Métricas}
    \begin{block}{Perda de Pacotes (\textit{Packet Loss})}
      \begin{itemize}
        \item Causada por congestionamento — roteadores descartam pacotes
        \item QoS define \textbf{quais} pacotes descartar (prioridade)
        \item Em tempo real: causa jitter e travamentos de fala/vídeo
      \end{itemize}
    \end{block}
    \vspace{0.5em}
    \begin{block}{Outras Métricas}
      \begin{itemize}
        \item \textbf{MOS}: qualidade de áudio, escala 1--5
        \item \textbf{Throughput}: taxa de entrega bem-sucedida de dados
        \item \textbf{Error Rate}: frequência de pacotes corrompidos
      \end{itemize}
    \end{block}
  \end{frame}

  % =============================================
  \section{Técnicas de QoS}
  % =============================================

  \begin{frame}
    \frametitle{Técnicas de QoS}
    \begin{itemize}
      \item<1-> \textbf{Marcação de Tráfego}
        \begin{itemize}
          \item CoS: marca no cabeçalho do frame (camada 2)
          \item DSCP: marca no cabeçalho do pacote (camada 3)
        \end{itemize}
      \item<2-> \textbf{Filas (\textit{Queuing})}
        \begin{itemize}
          \item \textit{Priority queuing}: filas dedicadas por prioridade
          \item Reduz chances de descarte para tráfego crítico
        \end{itemize}
      \item<3-> \textbf{RSVP} (\textit{Resource Reservation Protocol})
        \begin{itemize}
          \item Protocolo de sinalização de camada de rede
          \item Reserva recursos antes de iniciar transmissão
        \end{itemize}
      \item<4-> \textbf{\textit{Traffic Shaping}}
        \begin{itemize}
          \item Limita taxa de saída, suaviza picos de tráfego
        \end{itemize}
    \end{itemize}
  \end{frame}

  % =============================================
  \section{QoS e Multimídia}
  % =============================================

  \begin{frame}
    \frametitle{QoS e Multimídia}
    Tráfego de áudio e vídeo é \alert{altamente sensível} a delay, jitter,
    packet loss e largura de banda.
    \vspace{0.5em}
    \begin{itemize}
      \item<2-> Prioriza VoIP, videochamadas e \textit{live stream}
      \item<3-> Controla latência e jitter para sincronização A/V
      \item<4-> Reduz perda de pacotes priorizando mídia em congestionamento
      \item<5-> Aloca maior largura de banda para fluxos multimídia
      \item<6-> Enfileira pacotes para compartilhamento eficiente da rede
    \end{itemize}
  \end{frame}

  % =============================================
  \section{Vantagens e Melhores Práticas}
  % =============================================

  \begin{frame}
    \frametitle{Vantagens do QoS}
    \begin{enumerate}
      \item<1-> Priorização de aplicações críticas
      \item<2-> Melhor gerenciamento de recursos de rede
      \item<3-> Melhora da experiência do usuário
      \item<4-> Gerenciamento de tráfego ponto a ponto
      \item<5-> Prevenção de perda de pacotes
      \item<6-> Redução de latência
    \end{enumerate}
  \end{frame}

  \begin{frame}
    \frametitle{Melhores Práticas}
    \begin{itemize}
      \item<1-> Não definir limites de banda muito baixos (evita descarte excessivo)
      \item<2-> Considerar distribuição de pacotes entre filas
      \item<3-> Garantias de banda apenas em serviços específicos
      \item<4-> Priorizar por \textbf{um único tipo} (serviço OU política de segurança)
      \item<5-> Minimizar complexidade da configuração
      \item<6-> Usar UDP para testes, não superalocar throughput
    \end{itemize}
  \end{frame}

  % =============================================
  \section*{Referências}
  % =============================================

  \begin{frame}[allowframebreaks]
    \frametitle{Referências}
    \bibliographystyle{alpha}
    \bibliography{bibliografia}
  \end{frame}

\end{document}
